\section{The kernel rate of the actual network}
Let $M=2^d$, and identify output functions with vectors in $\mathbb R^M$.
Define
\begin{equation}\label{eq:kernel}
 \Psi(t)=\frac2\pi\arcsin t,\qquad
 K(x,y)=\Psi^{\circ(L-1)}\!\left(\frac{1+x\cdot y}{d+1}\right).
\end{equation}
The positive definiteness of $K$ is proved in the next section.
For each Boolean function $g$, put
\begin{equation}\label{eq:energy}
 E_K(g)=\min_{g(x)z(x)\ge a\ \forall x}\frac12 z^TK^{-1}z.
\end{equation}
The minimum exists since its feasible set is closed and nonempty and
the objective is coercive. Replacing the closed inequalities by strict
ones gives the same infimum, by scaling a feasible vector by $1+\epsilon$.

\begin{proposition}[Posterior concentration on minimum-energy signs]
\label{prop:rate}
For the actual network~\eqref{eq:weights}, at every fixed
$0<\gamma<1$,
\begin{equation}\label{eq:rate}
 \lim_{N\to\infty}N^{-\gamma}\log
 \mathbb P\{g(x)z_N(x)>a\ \forall x\}=-E_K(g).
\end{equation}
Consequently, after any fixed noiseless training sample, posterior
mass tends to one on the compatible functions minimizing $E_K$.
No assertion about the relative weights of tied minimizers is needed.
\end{proposition}

\begin{proof}
For hidden layer $s$, let
$Q_N^s(x,y)=N^{-1}\sum_i h_i^s(x)h_i^s(y)$.
We first establish, for each fixed $\epsilon>0$,
\begin{equation}\label{eq:freeze}
 \mathbb P\{\|Q_N^{L-1}-K\|>\epsilon\}
       \le C_\epsilon e^{-c_\epsilon N}
\end{equation}
for all sufficiently large $N$. Any matrix norm can be used, since
$M$ is fixed. Conditional on the preceding layer, the rows of the
next hidden layer are independent bounded vectors. Hoeffding's
inequality controls all their empirical second moments. Their
preactivations have covariance $N^{1-\gamma}Q_N^{s-1}$; at the first
layer the covariance is
$N^{1-\gamma}(1+x\cdot y)/d$. Because $N^{1-\gamma}\to\infty$,
the conditional second-moment map converges to the Gaussian-sign map
\[
 Q\longmapsto \left[\Psi\!\left(
 \frac{Q(x,y)}{\sqrt{Q(x,x)Q(y,y)}}\right)\right]_{x,y}.
\]
This follows by dominated convergence for a Gaussian vector, since
each scalar variance stays bounded away from zero in a neighborhood
of the limiting covariance. The convergence is uniform in such a
neighborhood: a compactness subsequence and a continuous square root
coupling reduce the assertion to the same dominated-convergence
argument. The limiting map is continuous there. Inducting through
the fixed number of layers, with smaller fixed tolerances at earlier
layers, proves~\eqref{eq:freeze}. The first-layer limit normalizes
the diagonal $(d+1)/d$, giving the kernel in~\eqref{eq:kernel}.

Conditionally on the hidden layers, the readout is exactly Gaussian:
\begin{equation}\label{eq:readout}
 Z_N\mid Q_N^{L-1}\sim\mathcal N(0,N^{-\gamma}Q_N^{L-1}).
\end{equation}
Couple it to $Y_N=N^{-\gamma/2}K^{1/2}\xi$, where
$\xi\sim\mathcal N(0,I_M)$ is independent of the hidden layers.
On a fixed sufficiently small neighborhood of the positive-definite
matrix $K$, the square-root map is Lipschitz. Hence, for $\delta>0$,
\[
 \mathbb P\{\|Z_N-Y_N\|>\delta\}
 \le C_\epsilon e^{-c_\epsilon N}
 +\mathbb P\{\|\xi\|>\delta N^{\gamma/2}/(C\epsilon)\}.
\]
First divide logarithms by $N^\gamma$ and let $N\to\infty$;
then let $\epsilon\downarrow0$. The resulting upper limit is
$-\infty$, since $\gamma<1$. The coupled vectors are therefore
exponentially equivalent at speed $N^\gamma$.

The fixed-dimensional Gaussian vectors $Y_N$ have large-deviation
rate $I(z)=z^TK^{-1}z/2$, directly from their densities. Exponential
equivalence transfers the same bounds to $Z_N$. Applying the upper
bound to a closed margin orthant and the lower bound to its interior
gives~\eqref{eq:rate}, since their infima agree. Finally, there are
only finitely many Boolean functions for fixed $d$. Dividing each
compatible orthant probability by their sum proves the posterior
concentration claim.
\end{proof}

This is a rate statement for the network itself. It does not require
relative-error approximation of rare probabilities, nor convergence
in total variation to a Gaussian posterior.

\begin{proposition}[The nonpositive-exponent posterior limits]
\label{prop:gp-limits}
For $\gamma=0$, the Boolean posterior converges to that of
$G\sim\mathcal N(0,K)$ conditioned on $|G(x)|>a$ everywhere and on
the observed signs. For $\gamma<0$, it converges to the Gaussian
posterior conditioned only on the observed signs.
\end{proposition}
\begin{proof}
The hidden-kernel argument establishing~\eqref{eq:freeze} holds
for every $\gamma<1$. By~\eqref{eq:readout}, the normalized vector
$N^{\gamma/2}Z_N$ therefore converges in law to
$\mathcal N(0,K)$. Its all-input margin is $aN^{\gamma/2}$, which
equals $a$ when $\gamma=0$ and tends to zero when $\gamma<0$.
The limiting Gaussian is nonsingular. Thus all boundaries involved
in any fixed full-sign event or margin event have probability zero,
and the conditioning events have positive probability. Taking
ratios of their probabilities gives convergence of each Boolean
posterior mass. There are finitely many masses for fixed $d$.
\end{proof}

